\section{程序结构与代码逻辑（伪代码形式呈现）}

\subsection{求解器结构体}

程序使用结构体统一保存：

\begin{itemize}[leftmargin=2em]
    \item 网格、时间和气体参数；
    \item 左右初始状态；
    \item 当前、预测和下一时间层守恒量；
    \item 当前与预测状态的通量；
    \item 原始量、声速及人工粘性传感器；
    \item 内存分配、初始化、推进、输出和释放等方法函数指针。
\end{itemize}

\begin{lstlisting}[caption={求解器结构体示意}]
typedef struct Riemann_1D_MacC_solver solver;

struct Riemann_1D_MacC_solver {
    int nx;
    double dx, dt, t, t_max, cfl, gamma;

    double *q1, *q2, *q3;
    double *q1_bar, *q2_bar, *q3_bar;
    double *q1_next, *q2_next, *q3_next;

    double *rho, *u, *p, *E, *a;

    int  (*allocate)(solver *self);
    void (*step_maccormack)(solver *self);
    void (*destroy)(solver *self);
};
\end{lstlisting}

\subsection{内存组织}

说明方案 B2.0：内存中只保留当前状态、预测状态和下一状态，计算完成后
用下一状态覆盖当前状态；中间过程通过定期写出快照文件保存。

\subsection{方法绑定与伪 OOP}

\begin{lstlisting}[caption={方法函数绑定示意}]
void solver_bind_methods(solver *self) {
    self->allocate = solver_allocate;
    self->step_maccormack = solver_step_maccormack;
    self->destroy = solver_destroy;
}
\end{lstlisting}

解释函数指针、\texttt{solver *self} 和
\texttt{self->member} 的含义。

\subsection{主程序流程}

\begin{enumerate}[leftmargin=2em]
    \item 创建默认求解器并绑定方法；
    \item 读取算例、网格和数值参数；
    \item 为所有数组分配堆内存；
    \item 初始化 Riemann 左右状态；
    \item 依据 CFL 计算时间步；
    \item 执行 MacCormack 预测与校正；
    \item 更新守恒量和原始量；
    \item 按间隔输出快照，最终输出结果；
    \item 释放动态内存。
\end{enumerate}
